% LaTeX Canvas: Chapter Outline
\chapter{Command Line}
\label{ch:terminal}

\section*{Purpose}
The terminal (command line) is a fast, precise interface for navigating files, running programs, and automating repetitive work. For novices, it can also be intimidating and risky because small mistakes can delete or overwrite work. This chapter builds confident, safe habits for everyday command-line use.

\section*{Learning objectives}
By the end of this chapter, you should be able to:
\begin{enumerate}
\item Explain what a \textbf{terminal} and a \textbf{shell} are, and how commands run.
\item Navigate the file system reliably (absolute/relative paths, current directory).
\item Create, view, copy, move, and delete files and directories safely.
\item Use help tools (\texttt{--help}, \texttt{man}) and interpret usage/syntax.
\item Combine commands with pipes and redirection to build small workflows.
\item Understand exit codes, errors, and common failure messages.
\item Use permissions and \texttt{sudo} responsibly (least privilege, when \emph{not} to use it).
\item Apply basic security hygiene: protect secrets, avoid unsafe pastes, and reduce accidental damage.
\end{enumerate}

\section*{Scope note}
This chapter focuses on Unix-like shells (macOS Terminal, Linux, and Windows via WSL). An appendix maps concepts to Windows PowerShell where relevant.

\section{A beginner mental model}
\subsection{Terminal vs shell}
\begin{itemize}
\item \textbf{Terminal:} the application window you type in.
\item \textbf{Shell:} the command interpreter (often \texttt{bash} or \texttt{zsh}) that runs programs.
\item \textbf{Command:} a program + arguments you ask the shell to run.
\end{itemize}

\subsection{What happens when you press Enter}
\begin{enumerate}
\item The shell reads your command line.
\item It expands shortcuts (globs, variables).
\item It locates the program (via \texttt{PATH}).
\item It runs the program in a process and prints output.
\item It returns an \textbf{exit code}: 0 usually means success.
\end{enumerate}

\subsection{Why the command line matters}
\begin{itemize}
\item High action-to-keystroke ratio (fast for repetitive work).
\item Composability (small tools + pipes).
\item Automation (scripts; reproducible workflows).
\item Remote access (SSH; see Chapter~\ref{ch:remote-computing}).
\end{itemize}

\section{Orientation and safety first}
\subsection{Know where you are (the current working directory)}
\begin{itemize}
\item Use \texttt{pwd} to print the current directory.
\item Use \texttt{ls} to list contents.
\item Use \texttt{cd} to change directories.
\end{itemize}

\subsection{The ``guardrails'' mindset}
\begin{itemize}
\item Prefer \textbf{read-only} commands first: \texttt{pwd}, \texttt{ls}, \texttt{cat}, \texttt{head}.
\item Before destructive actions: do a dry run (list the targets).
\item When in doubt: copy instead of move; move instead of delete.
\end{itemize}

\subsection{Common irreversible mistakes}
\begin{itemize}
\item Deleting the wrong directory (especially with recursive delete).
\item Running a command in the wrong directory.
\item Overwriting files with redirection (\texttt{>}).
\item Using \texttt{sudo} to ``fix'' permission problems without understanding them.
\end{itemize}

\section{Command anatomy and help}
\subsection{Basic syntax}
\begin{verbatim}
command [options] [arguments]
\end{verbatim}

\subsection{Flags and options}
\begin{itemize}
\item Short flags: \texttt{-l}, \texttt{-a}
\item Long flags: \texttt{--all}, \texttt{--help}
\item Some flags take values: \texttt{-o filename} or \texttt{--output=filename}
\end{itemize}

\subsection{Built-in help}
\begin{itemize}
\item \texttt{command --help}
\item \texttt{man command}
\item Skim \texttt{SYNOPSIS} first; then \texttt{OPTIONS}; then \texttt{EXAMPLES}.
\end{itemize}

\subsection{Quoting and spaces}
\begin{itemize}
\item Spaces separate arguments. Use quotes for filenames with spaces.
\item Single quotes usually prevent expansions; double quotes allow some expansions.
\item Rule: quote variables and paths unless you deliberately want splitting/globbing.
\end{itemize}

\section{File system navigation}
\subsection{Key path concepts}
\begin{itemize}
\item \texttt{/} root directory; \texttt{~} your home directory.
\item \texttt{.} current directory; \texttt{..} parent directory.
\item Absolute vs relative paths (Chapter~\ref{ch:filesystem} covers file system organization in depth).
\end{itemize}

\subsection{Core navigation commands}
\begin{itemize}
\item \texttt{pwd} (where am I?)
\item \texttt{ls} (what's here?) and useful flags (\texttt{-l}, \texttt{-a}, \texttt{-h})
\item \texttt{cd} (go there)
\end{itemize}

\subsection{Tab completion and history}
\begin{itemize}
\item Use \textbf{Tab} to complete commands and paths.
\item Use \textbf{Up Arrow} (and reverse search if taught) to reuse commands.
\item Best practice: edit a recalled command instead of retyping.
\end{itemize}

\section{Creating, viewing, and editing files}
\subsection{Creating directories and empty files}
\begin{itemize}
\item \texttt{mkdir} for directories (and \texttt{-p} for nested directories).
\item \texttt{touch} to create a file or update its timestamp.
\end{itemize}

\subsection{Viewing file contents safely}
\begin{itemize}
\item Small files: \texttt{cat}
\item Large files: \texttt{less} (scroll/search)
\item Previews: \texttt{head}, \texttt{tail}
\end{itemize}

\subsection{Creating small files from the command line}
\begin{itemize}
\item Redirection to create simple text files.
\item Prefer a text editor for multi-line content.
\end{itemize}

\subsection{Editors in the terminal}
\begin{itemize}
\item Use a beginner-friendly editor flow (e.g., \texttt{nano}) for quick edits.
\item Know how to exit without panic.
\item For longer-term work, use a full editor (VS Code, etc.).
\end{itemize}

\section{File operations: copy, move, delete (with safety)}
\subsection{Copy and move}
\begin{itemize}
\item \texttt{cp source dest} (copy)
\item \texttt{mv source dest} (move/rename)
\item Use interactive flags where available to prevent overwrites.
\end{itemize}

\subsection{Deleting files and directories}
\begin{itemize}
\item \texttt{rm file} deletes a file.
\item \texttt{rm -r dir} deletes a directory tree.
\item \textbf{High-risk:} \texttt{rm -rf}. Teach why it is dangerous.
\item Safe practice: \texttt{ls} the targets first; consider \texttt{rm -i}.
\end{itemize}

\subsection{Globbing (wildcards) and why it matters}
\begin{itemize}
\item \texttt{*} matches many characters; \texttt{?} matches one.
\item Globs expand \emph{before} the command runs.
\item Safety: run \texttt{echo *.csv} to preview what a glob matches.
\end{itemize}

\section{Searching and inspecting}
\subsection{Finding files}
\begin{itemize}
\item \texttt{find} for locating files by name, type, time.
\item Use carefully; start in a narrow directory scope.
\end{itemize}

\subsection{Searching within files}
\begin{itemize}
\item \texttt{grep} for finding patterns in text.
\item Teach minimal flags: recursive search, case-insensitive search.
\end{itemize}

\subsection{Inspecting file metadata}
\begin{itemize}
\item \texttt{ls -l} for permissions, size, timestamps.
\item \texttt{file} to identify file types.
\item Checksums concept (optional): how to compare transfers.
\end{itemize}

\section{Pipes and redirection: building workflows}
\subsection{Standard streams}
\begin{itemize}
\item \texttt{stdin} (input), \texttt{stdout} (output), \texttt{stderr} (errors).
\end{itemize}

\subsection{Redirection}
\begin{itemize}
\item \texttt{>} overwrite output to a file; \texttt{>>} append.
\item Safety: avoid accidental overwrite; teach creating backups.
\end{itemize}

\subsection{Pipes}
\begin{itemize}
\item \texttt{|} sends output of one command into the next.
\item Common patterns: \texttt{ls | less}, \texttt{grep pattern file | head}.
\end{itemize}

\subsection{Exit codes and ``did it work?''}
\begin{itemize}
\item Teach the idea: commands signal success/failure.
\item Encourage checking outputs, not assuming.
\end{itemize}

\section{Environment basics: \texttt{PATH}, variables, and reproducibility}
\subsection{Environment variables}
\begin{itemize}
\item What they are and why tools use them.
\item Basic operations: \texttt{echo \$VAR}, setting variables (conceptual).
\end{itemize}

\subsection{PATH and locating commands}
\begin{itemize}
\item \texttt{PATH} controls where the shell searches for programs.
\item Use \texttt{which} (or equivalent) to see which program will run.
\item Common pitfall: multiple versions of a tool installed.
\end{itemize}

\subsection{Working directory as a dependency}
\begin{itemize}
\item Relative paths depend on where you run the command.
\item Best practice: adopt a stable project root and use relative paths within it.
\end{itemize}

\section{Permissions, ownership, and \texttt{sudo}}
\subsection{The permission model (just enough)}
\begin{itemize}
\item Users, groups, and ``other.''
\item Read/write/execute bits and what they mean for files vs directories.
\item Recognizing permission errors.
\end{itemize}

\subsection{When to use \texttt{sudo} (and when not to)}
\begin{itemize}
\item \texttt{sudo} runs a single command with elevated privileges.
\item Use it for system-level changes (installing software, editing system configs) when appropriate.
\item Do \textbf{not} use \texttt{sudo} to fix a broken project folder or to run routine scripts.
\end{itemize}

\subsection{Safe \texttt{sudo} practices}
\begin{itemize}
\item Read the command twice before running.
\item Prefer explicit paths over wildcards.
\item Avoid piping into \texttt{sudo} unless you understand the implications.
\item Never run unreviewed copy-pasted commands with \texttt{sudo}.
\end{itemize}

\subsection{Recovering from permission issues}
\begin{itemize}
\item First identify ownership/permissions and the target directory.
\item If you created a mess with \texttt{sudo}, stop and ask for help.
\end{itemize}

\section{Security hygiene for terminal users}
\subsection{Secrets and command history}
\begin{itemize}
\item Your shell history may record commands.
\item Never put passwords/tokens directly on the command line if avoidable.
\item Prefer environment variables, config files with correct permissions, or secret managers (conceptual).
\end{itemize}

\subsection{Copy/paste safety}
\begin{itemize}
\item Understand that copied commands may include hidden characters.
\item Always inspect commands before running; avoid running multi-line scripts you do not understand.
\item Prefer official documentation sources.
\end{itemize}

\subsection{Least privilege}
\begin{itemize}
\item Work in your home directory as a normal user.
\item Elevate privileges only for specific administrative tasks.
\end{itemize}

\section{Workflow patterns for students}
\subsection{Pattern 1: ``enter a project, run, exit''}
\begin{enumerate}
\item \texttt{cd} into a project folder.
\item List files; confirm you are in the right place.
\item Run a command (script, notebook server, tests).
\item Save outputs into a clear directory.
\item Exit cleanly.
\end{enumerate}

\subsection{Pattern 2: creating a reproducible command log}
\begin{itemize}
\item Keep a \texttt{README.md} or \texttt{notes.txt} with the commands needed to reproduce results.
\item Record inputs, outputs, and environment assumptions.
\end{itemize}

\subsection{Pattern 3: safe cleanup}
\begin{itemize}
\item Identify targets with search and preview.
\item Archive important outputs.
\item Delete with care; verify afterward.
\end{itemize}

\section{Troubleshooting playbook}
\subsection{Common errors and what they mean}
\begin{description}
\item[\texttt{command not found}] Program not installed or PATH misconfigured.
\item[\texttt{No such file or directory}] Wrong path, wrong working directory, or typo.
\item[\texttt{Permission denied}] Insufficient rights; avoid reflexive \texttt{sudo}.
\item[\texttt{Is a directory}] Command expected a file.
\end{description}

\subsection{A disciplined response}
\begin{enumerate}
\item Re-run the command slowly; check spelling.
\item Print current directory; list files.
\item Confirm paths and file extensions.
\item Consult \texttt{--help} or \texttt{man}.
\item If still stuck: create a minimal reproduction and ask a precise question.
\end{enumerate}

\section{Worked examples (outline)}
\subsection{Example 1: Build a course workspace}
\begin{itemize}
\item Create directories for \texttt{courses/}, \texttt{assignments/}, \texttt{project/}.
\item Practice navigation and listing.
\end{itemize}

\subsection{Example 2: Create and inspect a dataset file}
\begin{itemize}
\item Download/move a \texttt{.csv}.
\item Preview with \texttt{head} and \texttt{less}.
\item Search for a column name with \texttt{grep}.
\end{itemize}

\subsection{Example 3: Build a small pipeline}
\begin{itemize}
\item Combine commands with pipes.
\item Redirect output to a file.
\item Verify results.
\end{itemize}

\subsection{Example 4: Diagnose and fix a ``file not found'' error}
\begin{itemize}
\item Identify working directory mismatch.
\item Repair using relative paths and project root conventions.
\end{itemize}

\section{Templates}
\subsection{Template A: Safe destructive action checklist}
\begin{verbatim}
Before deleting/moving:

1. pwd
2. ls (confirm context)
3. Preview targets (echo glob / ls targets)
4. Copy/backup if uncertain
5. Execute with least privilege (no sudo unless required)
6. Verify after
   \end{verbatim}

\subsection{Template B: Reproducible command log}
\begin{verbatim}
Project:
Date:
Goal:
Commands:

* ...
  Inputs:
* ...
  Outputs:
* ...
  Notes/assumptions:
* ...
  \end{verbatim}

\section{Exercises}
\begin{enumerate}
\item Navigate from your home directory to a course folder using only \texttt{cd}, \texttt{pwd}, \texttt{ls}.
\item Create a project directory structure with \texttt{mkdir -p}.
\item Create a text file, preview it with \texttt{head}, then open it in \texttt{less}.
\item Copy a directory; rename the copy; explain the difference between copy and move.
\item Use a wildcard to list only \texttt{.csv} files; preview the expansion with \texttt{echo}.
\item Build a pipeline that searches within a file and writes matching lines to an output file.
\item Trigger a permission error intentionally (in a safe location) and explain why \texttt{sudo} is not the first fix.
\end{enumerate}

\section{One-page checklist}
\begin{itemize}
\item I can explain terminal vs shell, and command + arguments.
\item I can navigate with \texttt{pwd}, \texttt{ls}, \texttt{cd} and understand paths.
\item I can create files/directories and view contents safely.
\item I can copy/move/delete with caution and preview wildcards.
\item I can use \texttt{--help} and \texttt{man}.
\item I can use pipes and redirection and understand overwrite vs append.
\item I understand permissions and use \texttt{sudo} only when appropriate.
\item I avoid leaking secrets in commands and history.
\end{itemize}

% \appendix
\section{Windows notes: PowerShell and WSL (optional)}
\subsection{Two common paths}
\begin{itemize}
\item \textbf{WSL} gives a Linux-like shell where \texttt{sudo} and Unix commands apply.
\item \textbf{PowerShell} uses different command names and syntax (but the same mental models: paths, pipes, permissions).
\end{itemize}

\subsection{Concept mapping (high level)}
\begin{itemize}
\item Navigation: \texttt{cd}, \texttt{pwd} (\texttt{Get-Location}), listing (\texttt{ls}/\texttt{dir}).
\item Help: \texttt{--help}/\texttt{man} vs \texttt{Get-Help}.
\item Pipes exist in both, but objects vs text differ.
\end{itemize}


\section{What is a terminal?}

A terminal is a tool for interacting with your computer by issuing commands. A terminal launches a shell (e.g.\ \texttt{bash} or \texttt{zsh}), so sometimes you will hear a ``terminal'' also described as a ``shell''.

\section{What terminal commands should I know, as a basic foundation?}

The exact syntax of terminal commands will be slightly different on Windows and Mac. So to keep things consistent, we are only going to detail high-level descriptions of terminal commands in this guide. 
You can Google the exact commands for your specific system; it is good practice to learn to search for the terminal commands you will need online.

To get started with the terminal for your INFO classes, you should memorize the commands for each of the following: 

\begin{itemize}
\item Activate/deactivate a conda environment
\item Change your directory
\item Install a specific version of a package with conda
\item Inspect the location of a program to the terminal (e.g.\ \texttt{which Python})
\item Print your working directory
\end{itemize}